\section{The finite outer-parallel volume bound}
\label{sec:global}

For $r\ge1$, define
\begin{equation}\label{eq:Ur}
 U_r=\bigcup_{i=1}^n(c_i+rB),
\end{equation}
where $B$ is the closed unit ball.  In the terminology of
\cite{BezdekLangi2025}, this is the outer-parallel domain of the hard unit-ball
packing with outer radius
\[
 \lambda=r-1.
\]

For an arbitrary unit-ball packing $\mathcal P$ with hard union $P$ and
outer-parallel union $P_\lambda$, Bezdek and L\'angi define the truncated density
\begin{equation}\label{eq:truncdef}
 \delta(\lambda,\mathcal P)
 =\limsup_{R\to\infty}
 \frac{\vol(P\cap RB)}{\vol(P_\lambda\cap RB)},
\end{equation}
and define $\delta(\lambda)$ as the supremum over all packings.  Their class of
packings includes finite packings.  For the present finite packing, once $R$ is
large enough to contain $P$ and $P_\lambda$, the quotient in
\cref{eq:truncdef} stabilizes.  The hard balls have pairwise disjoint interiors,
so their union has volume $n\vol(B)$, and therefore
\begin{equation}\label{eq:finitequotient}
 \delta(\lambda,\mathcal P)
 =\frac{n\vol(B)}{\vol(U_r)}.
\end{equation}

We use the following identity as a named published input.  Equation (1.2) of
\cite{BezdekLangi2025} states that
\begin{equation}\label{eq:BLidentity}
 \delta(\lambda)=\delta_3\qquad(\lambda\ge1),
\end{equation}
where $\delta_3$ is the ordinary packing density.  The finite specialization
\cref{eq:finitequotient} is the only additional reduction needed here.  The
Kepler theorem gives~\cite{Hales2005}
\[
 \delta_3=\frac{\pi}{\sqrt{18}}.
\]
At the optimized radius, \cref{eq:rstar} and the calculation in
\cref{sec:optimizer} give $r_*>2$, hence $\lambda_*=r_*-1>1$.  Since every
individual truncated density is at most its supremum,
\[
 \frac{n(4\pi/3)}{\vol(U_{r_*})}
 \le\frac{\pi}{\sqrt{18}},
\]
and therefore
\begin{equation}\label{eq:volumelower}
 \boxed{\vol(U_{r_*})\ge4\sqrt2\,n.}
\end{equation}
This is a finite statement for every $n\ge2$, conditional on the published
identity \eqref{eq:BLidentity}.

The union $U_{r_*}$ is bounded and has finite perimeter.  Euclidean
isoperimetry gives
\[
 \Per(U_{r_*})^3\ge36\pi\vol(U_{r_*})^2,
\]
and $\Per(U_{r_*})\le\cH^2(\partial U_{r_*})$.  Using
\cref{eq:volumelower},
\begin{equation}\label{eq:surfacelower}
 \boxed{
 \cH^2(\partial U_{r_*})
 \ge4\pi\left(\frac{18}{\pi^2}\right)^{1/3}n^{2/3}.}
\end{equation}

\section{Local spherical-neighborhood estimate}
\label{sec:local}

Fix a center $c_i$.  For each contact neighbor $c_j$, put
\[
 v_{ij}=\frac{c_j-c_i}{2}\in\Sph^2.
\]
If $j\ne k$ are contact neighbors, then $|c_j-c_k|\ge2$, hence
$|v_{ij}-v_{ik}|\ge1$, so their angular separation is at least $\pi/3$.
Consequently the interiors of the caps $C(v_{ij},\pi/6)$ are pairwise disjoint.

Set
\begin{equation}\label{eq:q}
 \beta=\frac\pi6,
 \qquad
 q=1-\cos\beta=1-\frac{\sqrt3}{2},
\end{equation}
and define the base set
\[
 A_i=\bigcup_{j\sim i} C(v_{ij},\beta).
\]
This is a finite union of closed caps and hence is compact and measurable.  Cap
boundaries are circles and have zero spherical area.  Therefore, if $d_i=d$,
\begin{equation}\label{eq:basearea}
 \area(A_i)=2\pi dq.
\end{equation}
Let $\rho_d$ be the radius of an equal-area cap.  Then
\begin{equation}\label{eq:xd}
 x_d:=\cos\rho_d=1-dq.
\end{equation}

A point $c_i+ru$ on the sphere of radius $r$ lies in the enlarged ball around a
contact neighbor in direction $v_{ij}$ exactly when
\[
 |ru-2v_{ij}|\le r,
\]
which is equivalent to $u\cdot v_{ij}\ge1/r$.  Thus the closed hidden cap has
angular radius
\begin{equation}\label{eq:alpha}
 \alpha=\arccos(1/r).
\end{equation}
For $r\ge2$, put $s=\alpha-\beta\ge\pi/6>0$.  Distance to a cap satisfies
\[
 \operatorname{dist}_{\Sph^2}(u,C(v,\beta))
 =\max\{0,\angle(u,v)-\beta\},
\]
so the union of the closed hidden caps is exactly the closed $s$-neighborhood of
$A_i$.

Spherical-neighborhood isoperimetry says that among measurable spherical sets of
given area, a cap minimizes every closed geodesic-neighborhood area
\cite{BaernsteinTaylor1976,BarnesOzgurWu2018}.  At the optimized radius and for
$d\le11$, non-saturation holds exactly.  The endpoint identity below gives
$\rho_{11}=\pi-\alpha_*$, and hence
\begin{equation}\label{eq:nonsaturation}
 \rho_d+s_*
 \le\rho_{11}+\alpha_*-\beta
 =\pi-\beta
 =\frac{5\pi}{6}<\pi.
\end{equation}
Therefore the exposed area $e_i$ on the sphere of radius $r_*$ satisfies
\begin{equation}\label{eq:exposure}
 e_i\le2\pi r_*^2H_d,
 \qquad
 H_d=1+\cos(\rho_d+s_*).
\end{equation}
The comparison above uses closed hidden caps, whereas the actually exposed patch
removes interiors of enlarged balls.  Their difference is contained in a finite
union of sphere-intersection circles and therefore has $\cH^2$-measure zero.
Non-tangent enlarged balls can only reduce actual exposure.
